\chapter{Valorisation par transformée de Fourier et ingénierie de la calibration}\label{ch:dv:fourier-pricing-and-calibration-engineering}

Chaque matin, un job ajuste un \omterm{def:dv:stochastic-volatility:heston}{modèle de Heston} à vingt cotations d'options. Au jour 58 du marché synthétique de ce chapitre, la \omterm{def:dv:stochastic-volatility:sv}{volatilité
de la volatilité} ajustée vaut 0.51~; au jour 59, elle vaut 0.28, l'erreur d'ajustement restant inchangée à 0.31 point de volatilité. Rien ne s'est produit sur le marché que les
cotations puissent voir~: le mouvement relève du bruit, amplifié par une direction de l'espace des paramètres le long de laquelle l'ajustement varie à peine. Mais les produits exotiques valorisés
avec le modèle les voient. L'ajustement de convexité d'un \omterm{def:dv:variance-swaps-and-volatility-derivatives:volswap}{swap de volatilité} à un an passe de 1.32 à 0.39 point de volatilité en une nuit, et un
desk qui valorise son livre de cette façon constate un profit ou une perte qu'aucune couverture ne saurait expliquer. Ce chapitre traite des deux moitiés du travail~: valoriser
assez vite à partir d'une fonction caractéristique pour calibrer dessus, et calibrer d'une manière dont le résultat peut servir à un desk d'un
jour sur l'autre.

\section{Valorisation par transformée}

Les chapitres~10 et~13 ont écrit chaque modèle à travers la fonction caractéristique $\varphi(u)=\E[e^{iux_T}]$ de $x_T=\ln(S_T/F)$, et ont
valorisé au moyen d'une seule intégrale. Cette intégrale découle de l'identité de Parseval.

\begin{definition}[Formule de Lewis]\label{def:dv:fourier-pricing-and-calibration-engineering:lewis}
La \emph{formule de Lewis}\index{formule de Lewis} valorise un call européen à partir de la fonction caractéristique de $x_T$ le long de la
droite $\operatorname{Im}u=-\frac12$~: avec $k=\ln(K/F)$,
\[
C=P(0,T)\Bigl(F-\frac{\sqrt{FK}}{\pi}\int_0^\infty\operatorname{Re}\bigl(e^{-iuk}\varphi(u-\tfrac i2)\bigr)\frac{du}{u^2+\frac14}\Bigr).
\]
\end{definition}

\begin{proposition}[Démonstration]\label{prop:dv:fourier-pricing-and-calibration-engineering:lewis}
Avec $g(x)=\min(e^x,e^k)$, on a $F-C/P(0,T)=F\,\E[g(x_T)]$. Pour $0<\operatorname{Im}z<1$, la transformée $\hat g(z)=\int e^{izx}g(x)\,dx$ vaut
\[
\hat g(z)=\int_{-\infty}^k e^{(1+iz)x}dx+e^k\int_k^\infty e^{izx}dx=\frac{e^{(1+iz)k}}{1+iz}-\frac{e^{(1+iz)k}}{iz}=\frac{e^{(1+iz)k}}{z^2-iz},
\]
et l'identité de Parseval donne $\E[g(x_T)]=\frac1{2\pi}\int\varphi(-z)\hat g(z)\,du$ le long de $z=u+\frac i2$. Là, $z^2-iz=u^2+\frac14$ et
$e^{(1+iz)k}=e^{k/2}e^{iuk}$. En remplaçant $u$ par $-u$ puis en repliant l'intégrale sur $u>0$, on obtient la formule, puisque $Fe^{k/2}=\sqrt{FK}$.
\end{proposition}

Le même raisonnement valorise tout payoff dont la transformée est connue sur une bande, c'est là le point de Lewis. Une seconde voie amortit le call lui-même.

\begin{definition}[Formule de Carr--Madan]\label{def:dv:fourier-pricing-and-calibration-engineering:cm}
La \emph{formule de Carr--Madan}\index{formule de Carr--Madan} écrit le call amorti $e^{\alpha k}C(k)$, $\alpha>0$, comme une intégrale de
Fourier en le log-strike~: avec $F=1$,
\[
C(k)=\frac{e^{-\alpha k}}{\pi}\int_0^\infty\operatorname{Re}\bigl(e^{-ivk}\psi(v)\bigr)dv,\qquad
\psi(v)=\frac{P(0,T)\,\varphi\bigl(v-(\alpha+1)i\bigr)}{\alpha^2+\alpha-v^2+i(2\alpha+1)v},
\]
qui, sur une grille en $v$ et en $k$, est une transformée de Fourier discrète~: la transformée de Fourier rapide renvoie les calls sur
toute une grille de log-strikes d'un seul coup.
\end{definition}

La méthode COS (livre~4, chapitre~28) emprunte une troisième voie. Elle développe la densité de $\ln(S_T/K)$ en série de cosinus sur un
intervalle tronqué $[a,b]$. Les coefficients proviennent directement de $\varphi$, et le payoff du put s'intègre contre chaque cosinus sous forme fermée.
Pour les densités régulières, l'erreur décroît exponentiellement avec le nombre de termes.

\begin{example}[Trois pricers, un même smile]\label{ex:dv:fourier-pricing-and-calibration-engineering:cos}
Calls à un an dans le \omterm{def:dv:stochastic-volatility:heston}{modèle de Heston} de base du chapitre~10, sur 17 strikes de 60 à 140, face à une référence (l'intégrale de Lewis par
panneaux de Gauss--Legendre jusqu'à $u=1\,000$, précise à environ $10^{-9}$~; à la monnaie elle donne 6.7509). La plus grande erreur de la méthode COS vaut 0.012 avec 64 termes,
$8.6\times10^{-5}$ avec 128 et $4.1\times10^{-7}$ avec 256. Carr--Madan, avec interpolation entre les log-strikes de la grille, atteint
$2.1\times10^{-4}$ avec 1\,024 points, $1.8\times10^{-6}$ avec 4\,096, et stagne autour de $2\times10^{-7}$
(voir la \cref{fig:dv:fourier-pricing-and-calibration-engineering:transform}). L'intervalle de troncature compte autant que le nombre de termes. Avec $[a,b]$ à
douze écarts-types au lieu de seize, COS s'arrête à $3.0\times10^{-5}$ quel que soit le nombre de termes, car la queue gauche de
Heston est plus épaisse que celle d'une loi normale.
\end{example}

\begin{omfigure}
\begin{tikzpicture}
\begin{axis}[width=10.5cm,height=5cm,xmode=log,ymode=log,log basis x=2,
  xlabel={termes (COS) ou points de grille (Carr--Madan)},ylabel={erreur de prix maximale},
  tick label style={font=\small},label style={font=\small},
  xmin=12,xmax=22000,ymin=1e-7,ymax=3,grid=major,grid style={gray!25},
  xtick={16,64,256,1024,4096,16384},xticklabels={16,64,256,1\,024,4\,096,16\,384},
  legend columns=2,legend cell align=left,legend style={font=\footnotesize,at={(0.5,-0.3)},anchor=north,draw=none,fill=none,/tikz/every even column/.append style={column sep=8pt}},
  table/col sep=comma]
\addplot[omThm,very thick,mark=*,mark size=1.6pt] table[x=n,y=err]{figdata/derivatives/24-fourier-pricing-and-calibration-engineering/cos.csv};
\addplot[omDef,very thick,mark=square*,mark size=1.6pt] table[x=n,y=err]{figdata/derivatives/24-fourier-pricing-and-calibration-engineering/fft.csv};
\legend{méthode COS,Carr--Madan FFT}
\end{axis}
\end{tikzpicture}

\omcaption{Plus grande erreur sur 17 strikes d'un call de Heston à un an, en fonction du nombre de termes ou de points de grille~: COS converge exponentiellement vers
la référence~; la FFT converge vite puis stagne à cause de ses erreurs d'interpolation et d'amortissement. Données~: le
tutoriel.}\label{fig:dv:fourier-pricing-and-calibration-engineering:transform}
\end{omfigure}

Pour la calibration, le choix est pratique. COS valorise les strikes qu'un desk cote, là où ils se trouvent, en quelques centaines d'opérations chacun. La FFT
valorise une grille fixe, et les strikes cotés demandent une interpolation. Un smile de cinq strikes coûte 800 évaluations de la fonction caractéristique avec COS à 160 termes, et 4\,096 avec la FFT.
Une calibration de Levenberg--Marquardt demande six évaluations du résidu par itération (cinq pour la jacobienne, une pour le pas) et, partant de
la veille, environ huit itérations~: soit quelque deux cents smiles.

\section{La calibration comme pipeline}

Le livre~4, chapitre~24, a traité la calibration comme un problème inverse. Sur un desk, c'est un job qui tourne chaque matin sur des données que personne n'a vérifiées. La
construction est un pipeline à cinq étapes, chacune avec ses propres modes de défaillance.
\begin{enumerate}
\item \textbf{Cotations.} Volatilités cotées au cours acheteur et au cours vendeur par échéance et par strike, avec le forward et la courbe d'actualisation. Un jeu de cotations est un instantané,
conservé afin que toute calibration passée puisse être rejouée.
\item \textbf{Filtre.} Écarter les cotations sans cours acheteur, les cotations croisées (cours acheteur supérieur au cours vendeur), les cotations périmées (fourchettes trop larges pour porter une information) et
les strikes trop éloignés pour compter. Le marché synthétique de ce chapitre corrompt chaque jour une cotation dans chaque sens. Le filtre écarte ces deux-là et conserve
les dix-huit autres.
\item \textbf{Objectif.} Erreurs pondérées et, éventuellement, une pénalité (sections suivantes).
\item \textbf{Optimiseur.} Levenberg--Marquardt, partant des paramètres de la veille, dans des coordonnées non contraintes
$z=(\ln v_0,\ln\kappa,\ln\bar v,\ln\eta,\operatorname{artanh}\rho)$ afin que chaque pas soit admissible.
\item \textbf{Diagnostics.} Erreur d'ajustement globale et par échéance, conditionnement de la jacobienne, variations des paramètres par rapport aux bornes. Une alarme
empêche les paramètres d'atteindre les pricers tant que quelqu'un ne les a pas examinés.
\end{enumerate}

\section{Objectif, pondérations et contraintes}

Calibrer en volatilité implicite est naturel, puisque c'est ainsi que les erreurs sont jugées. Mais chaque résidu demande alors une inversion de la
volatilité implicite. Calibrer en prix est plus rapide et exige un poids par cotation.

\begin{definition}[Pondération par le vega]\label{def:dv:fourier-pricing-and-calibration-engineering:vega}
La \emph{pondération par le vega}\index{pondération par le vega} divise chaque erreur de prix par le \omterm{def:dv:greeks-and-the-hedging-pnl:greeks}{vega} de Black de sa cotation, de sorte que $(C^{\text{modèle}}-C^{\text{milieu}})/
\mathcal V\approx\sigma^{\text{modèle}}-\sigma^{\text{milieu}}$~: une calibration en prix minimise alors les erreurs de volatilité implicite au premier
ordre, sans inversion.
\end{definition}

Sans pondération, les erreurs de prix sur les options à la monnaie de longue échéance dominent, puisque leurs \omterm{def:dv:greeks-and-the-hedging-pnl:greeks}{vegas} sont les plus grands.

\begin{example}[Ce qu'apportent les pondérations]\label{ex:dv:fourier-pricing-and-calibration-engineering:weights}
Le premier jour du marché synthétique (quatre échéances, d'un mois à un an, cinq strikes chacune, dix-huit après filtrage), l'erreur
quadratique moyenne de l'ajustement de Heston par échéance, en points de volatilité, est~:

\begin{center}\small
\begin{tabular}{lrrrrr}
\toprule
pondérations & 1 mois & 3 mois & 6 mois & 1 an & toutes\\
\midrule
\omterm{def:dv:greeks-and-the-hedging-pnl:greeks}{vega} & 0.35 & 0.21 & 0.14 & 0.11 & 0.22\\
égales (erreurs de prix) & 0.46 & 0.15 & 0.12 & 0.06 & 0.24\\
\bottomrule
\end{tabular}
\end{center}

Des pondérations égales donnent un meilleur ajustement à un an, au prix d'un moins bon ajustement à un mois. À la monnaie, le \omterm{def:dv:greeks-and-the-hedging-pnl:greeks}{vega} d'une option à un mois est égal à celui de l'option à un an divisé par
$\sqrt{12}$, de sorte qu'en prix ses erreurs comptent douze fois moins.
\end{example}

Les contraintes viennent du modèle et du desk. La transformation garde $v_0,\kappa,\bar v,\eta>0$ et $|\rho|<1$. La condition de Feller
n'est pas imposée. Elle n'est pas vérifiée dans la plupart des calibrations sur actions, celle-ci comprise ($2\kappa\bar v=0.19$ contre $\eta^2=0.24$), et les schémas de simulation
du chapitre 23 sont conçus pour ce cas.

\section{Stabilité d'un jour à l'autre}

Une calibration peut reproduire parfaitement les cotations et rester inutile, si une petite variation des cotations déplace fortement les paramètres. Cont et
Ben Hamida ont montré que l'erreur en échantillon peut présenter de nombreux minima presque globaux, de sorte que le problème de calibration est mal posé. L'erreur d'ajustement ne dit alors
rien sur le minimum retenu.

\begin{example}[Une vallée plate]\label{ex:dv:fourier-pricing-and-calibration-engineering:valley}
L'ajustement du premier jour donne $v_0=0.0451$, $\kappa=1.82$, $\bar v=0.0518$, $\eta=0.488$, $\rho=-0.632$, avec une erreur de 0.22 point de volatilité. En maintenant
$\eta$ fixé et en recalibrant les quatre autres, on obtient des erreurs de 0.29 pour $\eta=0.35$, 0.23 pour 0.45 et 0.55, et 0.27 pour 0.65
(la \cref{fig:dv:fourier-pricing-and-calibration-engineering:valley}, à gauche). Sur toute la plage, l'erreur varie de moins que le bruit des
cotations, soit 0.3 point de volatilité, tandis que $\kappa$ suit $\eta$. Les données de marché (un \omterm{def:dv:jumps-and-levy-models:bates}{modèle de Bates}, chapitre 13, avec bruit) ne permettent pas de distinguer ces
\omterm{def:dv:stochastic-volatility:heston}{modèles de Heston} entre eux.
\end{example}

\begin{omfigure}
\begin{tikzpicture}
\begin{groupplot}[group style={group size=2 by 1,horizontal sep=1.6cm},
  width=6.6cm,height=5cm,
  tick label style={font=\small},label style={font=\small},grid=major,grid style={gray!25},table/col sep=comma]
\nextgroupplot[xlabel={volatilité de la volatilité $\eta$ (fixée)},ylabel={erreur d'ajustement (points de volatilité)},xmin=0.2,xmax=1.05,ymin=0,ymax=0.45,
  title={\small profil, premier jour}]
\addplot[omThm,very thick,mark=*,mark size=1.6pt] table[x=eta,y=rmse]{figdata/derivatives/24-fourier-pricing-and-calibration-engineering/profile.csv};
\addplot[gray,dashed,domain=0.2:1.05] {0.3};
\nextgroupplot[xlabel={$\eta$ ajusté},ylabel={$\kappa$ ajusté},xmin=0.2,xmax=0.8,ymin=0.5,ymax=4.5,title={\small 60 ajustements quotidiens},
  legend columns=1,legend cell align=left,legend style={font=\footnotesize,at={(0.5,-0.32)},anchor=north,draw=none,fill=none}]
\addplot[omProp,only marks,mark=*,mark size=1.3pt] table[x=eta_free,y=kappa_free]{figdata/derivatives/24-fourier-pricing-and-calibration-engineering/daily.csv};
\addplot[omThm,only marks,mark=square*,mark size=1.3pt] table[x=eta_reg,y=kappa_reg]{figdata/derivatives/24-fourier-pricing-and-calibration-engineering/daily.csv};
\addplot[black,thick,dashed] table[x=eta,y=kappa]{figdata/derivatives/24-fourier-pricing-and-calibration-engineering/profile.csv};
\legend{sans pénalité,pénalité $10^{-3}$,vallée (profil)}
\end{groupplot}
\end{tikzpicture}

\omcaption{À gauche~: la meilleure erreur d'ajustement, la \omterm{def:dv:stochastic-volatility:sv}{volatilité de la volatilité} étant fixée, premier jour~; la ligne en pointillés figure le bruit des cotations. À droite~: les
couples $(\eta,\kappa)$ ajustés sur 60 jours, sans et avec une pénalité sur les variations des paramètres, et la vallée du premier jour (le meilleur $\kappa$ pour chaque
$\eta$ fixé)~: sans pénalité, les ajustements se dispersent largement~; avec pénalité, ils restent dans un amas serré.
Données~: le tutoriel.}\label{fig:dv:fourier-pricing-and-calibration-engineering:valley}
\end{omfigure}

Dériver le long de la vallée coûte de l'argent, car les produits exotiques ne se valorisent pas avec les seules vanilles ajustées.

\begin{definition}[P\&L de recalibration]\label{def:dv:fourier-pricing-and-calibration-engineering:recalib}
Le \emph{P\&L de recalibration}\index{P\&L de recalibration} d'une position est la variation de sa valeur de modèle provoquée par la recalibration des
paramètres, les entrées de marché sur lesquelles portent ses couvertures étant maintenues fixes~; dans une attribution du P\&L, il tombe dans la colonne
inexpliquée, et il n'a aucune couverture.
\end{definition}

Le remède consiste à indiquer à l'optimiseur que les paramètres de la veille sont informatifs. Avec $z$ les paramètres non contraints, la fonction objectif ajoute
$\lambda\lVert z-z_{\text{veille}}\rVert^2$ aux erreurs quadratiques pondérées, une pénalité de Tikhonov (livre~4, chapitre~24). Le long de la vallée, la
pénalité domine et les paramètres restent en place. En travers, ce sont les cotations qui dominent, et les paramètres bougent quand le marché bouge.

\begin{example}[Soixante jours avec et sans pénalité]\label{ex:dv:fourier-pricing-and-calibration-engineering:daily}
Sans pénalité, les ajustements quotidiens déplacent $\eta$ jusqu'à 0.235 en un jour, et l'ajustement de convexité d'un \omterm{def:dv:variance-swaps-and-volatility-derivatives:volswap}{swap de volatilité} à un an (chapitre 14)
varie de 0.36 point de volatilité par jour (écart type). Avec $\lambda=10^{-3}$, le plus grand déplacement quotidien de $\eta$ est de 0.071, et les
variations quotidiennes de l'ajustement tombent à 0.07 point de volatilité, et l'erreur d'ajustement moyenne passe de 0.291 à 0.303 point de volatilité
(la \cref{fig:dv:fourier-pricing-and-calibration-engineering:daily}). Le jour 59, où l'ajustement libre a fait passer $\eta$ de 0.51 à 0.28, l'ajustement
pénalisé se déplace de 0.18 point de volatilité au lieu de 0.93.
\end{example}

\begin{omfigure}
\begin{tikzpicture}
\begin{groupplot}[group style={group size=2 by 1,horizontal sep=1.6cm},
  width=6.6cm,height=5cm,xlabel={jour},xmin=0,xmax=61,
  tick label style={font=\small},label style={font=\small},grid=major,grid style={gray!25},table/col sep=comma]
\nextgroupplot[ylabel={volatilité de la volatilité $\eta$},ymin=0.2,ymax=0.8,
  legend columns=2,legend cell align=left,legend style={font=\footnotesize,at={(1.12,-0.3)},anchor=north,draw=none,fill=none,/tikz/every even column/.append style={column sep=8pt}}]
\addplot[omProp,thick] table[x=day,y=eta_free]{figdata/derivatives/24-fourier-pricing-and-calibration-engineering/daily.csv};
\addplot[omThm,very thick] table[x=day,y=eta_reg]{figdata/derivatives/24-fourier-pricing-and-calibration-engineering/daily.csv};
\legend{sans pénalité,pénalité $\lambda=10^{-3}$}
\nextgroupplot[ylabel={ajust. de convexité (pts de vol.)},ymin=0,ymax=2]
\addplot[omProp,thick] table[x=day,y=conv_free]{figdata/derivatives/24-fourier-pricing-and-calibration-engineering/daily.csv};
\addplot[omThm,very thick] table[x=day,y=conv_reg]{figdata/derivatives/24-fourier-pricing-and-calibration-engineering/daily.csv};
\end{groupplot}
\end{tikzpicture}

\omcaption{Soixante calibrations quotidiennes de Heston sur un marché synthétique bruité. À gauche~: la \omterm{def:dv:stochastic-volatility:sv}{volatilité de la volatilité} ajustée. À droite~: l'ajustement de
convexité d'un \omterm{def:dv:variance-swaps-and-volatility-derivatives:volswap}{swap de volatilité} à un an (volatilité du \omterm{def:dv:variance-swaps-and-volatility-derivatives:vs}{swap de variance} moins prix du \omterm{def:dv:variance-swaps-and-volatility-derivatives:volswap}{swap de volatilité}) implicite dans chaque
ajustement quotidien. Données~: le tutoriel.}\label{fig:dv:fourier-pricing-and-calibration-engineering:daily}
\end{omfigure}

Le poids de la pénalité est un choix. Trop faible, et les paramètres dérivent~; trop fort, et ils prennent du retard sur un marché qui bouge. La construction choisit
en rejouant l'historique~: pour chaque $\lambda$ d'une grille, on lance la calibration quotidienne sur les anciens jeux de cotations, on relève le plus grand déplacement quotidien des paramètres
et l'erreur d'ajustement moyenne, puis on retient le plus petit poids qui respecte les limites du desk
(la \cref{fig:dv:fourier-pricing-and-calibration-engineering:scan}).

\begin{omfigure}
\begin{tikzpicture}
\begin{groupplot}[group style={group size=2 by 1,horizontal sep=2.2cm},
  width=6.4cm,height=5cm,xmode=log,xlabel={poids de la pénalité $\lambda$},
  yticklabel style={/pgf/number format/fixed,/pgf/number format/precision=2},xmin=5e-6,xmax=2e-2,
  tick label style={font=\small},label style={font=\small},grid=major,grid style={gray!25},table/col sep=comma]
\nextgroupplot[ylabel={plus grand déplacement quotidien de $\eta$},ymin=0,ymax=0.25]
\addplot[omThm,very thick,mark=*,mark size=1.6pt] table[x=reg,y=max_d_eta]{figdata/derivatives/24-fourier-pricing-and-calibration-engineering/scan.csv};
\addplot[gray,dashed,domain=5e-6:2e-2] {0.1};
\nextgroupplot[ylabel={hausse de l'erreur d'ajust. (pts de vol.)},ymin=-0.01,ymax=0.12]
\addplot[omDef,very thick,mark=square*,mark size=1.6pt] table[x=reg,y=d_rmse]{figdata/derivatives/24-fourier-pricing-and-calibration-engineering/scan.csv};
\addplot[gray,dashed,domain=5e-6:2e-2] {0.1};
\end{groupplot}
\end{tikzpicture}

\omcaption{Rejeu de 60 jours pour chaque poids de pénalité. À gauche~: le plus grand déplacement quotidien de la \omterm{def:dv:stochastic-volatility:sv}{volatilité de la volatilité}, face à une limite de 0.1.
À droite~: la hausse de l'erreur d'ajustement moyenne, face à une limite de 0.1 point de volatilité. Données~: le
tutoriel.}\label{fig:dv:fourier-pricing-and-calibration-engineering:scan}
\end{omfigure}

\section{Tester une calibration}

Une calibration est du code et reçoit les tests de tout code, plus trois tests qui lui sont propres.
\begin{itemize}
\item \textbf{Aller-retour.} Générer des cotations à partir de paramètres connus et les retrouver. À partir de cotations produites par $v_0=0.045$, $\kappa=2$, $\bar v=0.05$,
$\eta=0.5$, $\rho=-0.65$, la chaîne retrouve chaque paramètre à $10^{-11}$ près, à partir d'un point de départ fourni par les paramètres de base du chapitre 10. Ce test éprouve le
pricer, la transformation et l'optimiseur ensemble.
\item \textbf{Identification.} Ajouter du bruit et recalibrer de nombreuses fois. Avec un bruit de 0.3 point de volatilité, vingt recalibrations donnent des écarts-types de
0.0007 pour $v_0$ (0.16 point de volatilité sur la volatilité à la monnaie), 0.049 pour $\eta$, 0.046 pour $\rho$, 0.0024 pour $\bar v$ et 0.56 pour
$\kappa$, soit 28~\% de sa valeur. Les cotations fixent le niveau de court terme et le skew, non la vitesse de retour à la moyenne. Le conditionnement de la
jacobienne, 30 ici, dit la même chose en un seul nombre.
\item \textbf{Rejeu.} Relancer les anciens jeux de cotations avec les réglages de production et vérifier les trajectoires des paramètres par rapport aux limites d'alarme, comme dans
la \cref{fig:dv:fourier-pricing-and-calibration-engineering:scan}.
\end{itemize}
Seul l'aller-retour admet une réponse juste. Les deux autres mesurent ce que les cotations permettent de soutenir, ce qui décide de la manière dont les exotiques dépendant de $\kappa$
doivent être réservés (chapitre 27).

\section{Tutoriel~: transformées et calibration quotidienne}\label{tut:dv:fourier-pricing-and-calibration-engineering:tut}

\begin{tutorial}
\textbf{Objectif.} Valoriser un smile de Heston par COS, Carr--Madan et la référence de Lewis~; calibrer un \omterm{def:dv:stochastic-volatility:heston}{modèle de Heston} chaque jour pendant 60 jours d'un marché synthétique bruité,
avec et sans pénalité sur les changements de paramètres, et choisir le poids de la pénalité par rejeu. \textbf{État final~:} les quatre figures,
les tableaux des poids et les nombres du problème du week-end.
\begin{tutsteps}
\item \textbf{Le pricer COS}, une ligne de termes en cosinus par strike~:
\omcode{code/firm/calib/firm_calib.py}{30}{48}{Calls par la méthode COS.}\label{lst:dv:fourier-pricing-and-calibration-engineering:cos}
\item \textbf{La calibration}~: résidus pondérés par le \omterm{def:dv:greeks-and-the-hedging-pnl:greeks}{vega}, pénalité sur la variation par rapport à la veille, Levenberg--Marquardt, diagnostics~:
\omcode{code/firm/calib/firm_calib.py}{164}{199}{Levenberg--Marquardt avec pénalité sur les variations de paramètres.}\label{lst:dv:fourier-pricing-and-calibration-engineering:calib}
\item \textbf{Exécuter} \texttt{dv\_calib.transform\_study()}, \texttt{weights\_table()}, \texttt{profile()}, \texttt{daily(0.0)}, \texttt{scan()},
\texttt{round\_trip()} et \texttt{fig\_calib.py}.
\end{tutsteps}
\textbf{Pour aller plus loin.} Calibrer un \omterm{def:dv:jumps-and-levy-models:bates}{modèle de Bates} sur les mêmes cotations et observer la dérive des paramètres de saut~; pénaliser seulement $\kappa$ et $\eta$~;
remplacer la pénalité de Tikhonov par un a priori centré sur une moyenne de long terme plutôt que sur la veille.
\end{tutorial}

\section{Construction~: la chaîne de calibration}\label{bld:dv:fourier-pricing-and-calibration-engineering:calib}

\begin{build}
\bfield{Rôle} Le service de calibration de la firme miniature~: pricers par transformée pour tout modèle à fonction caractéristique, et une chaîne quotidienne
qui transforme les cotations en paramètres sur lesquels les pricers peuvent compter.
\bfield{Interface} \texttt{cos\_calls(model, fwd, strikes, t, n, width)}, \texttt{carr\_madan(model, fwd, t, n, eta, alpha)},
\texttt{carr\_madan\_at}, \texttt{lewis\_calls} (référence), \texttt{cumulants}~; \texttt{QuoteSet(t, strike, bid, ask, fwd)},
\texttt{filter\_quotes(q, max\_spread, max\_sd)}~; \texttt{calibrate(q, z0, make, weights, prior, reg)} $\to$ \texttt{Fit(model, z, rmse,
penalty, iterations, condition)}~; \texttt{vol\_errors}, \texttt{jump\_alarm(z\_new, z\_old, limits)}~; \texttt{heston\_from}, \texttt{heston\_to}.
\bfield{Règles} Les jeux de cotations sont conservés et toute calibration peut être rejouée~; le filtre signale ce qu'il a écarté et pourquoi~; les poids sont explicites~;
la calibration part des paramètres de la veille~; une alarme bloque les paramètres qui sortent de leurs limites.
\bfield{Tests d'acceptation} \texttt{code/firm/calib/tests/}~: COS, Carr--Madan et la référence de Lewis, comparés à la formule de Black et entre eux~;
les cumulants d'une loi lognormale~; les quatre raisons du filtre~; l'aller-retour~; une pénalité infinie conserve les paramètres de la veille~; l'erreur de prix pondérée par le
\omterm{def:dv:greeks-and-the-hedging-pnl:greeks}{vega} est égale à l'erreur de volatilité au premier ordre.
\bfield{Pour aller plus loin} Le moteur FFT pour les modèles à dépendance temporelle du chapitre 12~; les jacobiennes analytiques issues des dérivées de la fonction caractéristique~;
une calibration du \omterm{def:dv:jumps-and-levy-models:bates}{modèle de Bates} avec les paramètres de saut pénalisés vers un a priori de long terme.
\end{build}

\begin{omsources}
\item P.~Carr et D.~B.~Madan, «~Option valuation using the fast Fourier transform~», \emph{Journal of Computational Finance} 2(4) (1999) 61--73.
\item A.~L.~Lewis, «~A simple option formula for general jump-diffusion and other exponential Lévy processes~», document de travail (2001).
\item F.~Fang et C.~W.~Oosterlee, «~A novel pricing method for European options based on Fourier-cosine series expansions~», \emph{SIAM Journal on
Scientific Computing} 31(2) (2008) 826--848.
\item R.~Cont et S.~Ben Hamida, «~Recovering volatility from option prices by evolutionary optimization~», \emph{Journal of Computational Finance}
8(4) (2005) 43--76.
\end{omsources}

\section{Exercices}

\begin{exercise}[$\star$]\label{exo:dv:fourier-pricing-and-calibration-engineering:1}
Vérifier que la \omterm{def:dv:fourier-pricing-and-calibration-engineering:lewis}{formule de Lewis} donne $C=P(0,T)F$ en $K=0$ et que $\varphi(-i)=1$ est ce qui fait de $x_T$ un log-prix forward-martingale.
\end{exercise}

\begin{exercise}[$\star$]\label{exo:dv:fourier-pricing-and-calibration-engineering:2}
Expliquer pourquoi la \omterm{def:dv:fourier-pricing-and-calibration-engineering:cm}{formule de Carr--Madan} a besoin du facteur d'amortissement $e^{\alpha k}$, et ce qui limite $\alpha$ par le haut.
\end{exercise}

\begin{exercise}[$\star$]\label{exo:dv:fourier-pricing-and-calibration-engineering:3}
Montrer qu'une erreur de prix divisée par le \omterm{def:dv:greeks-and-the-hedging-pnl:greeks}{vega} de Black est l'erreur de volatilité implicite au premier ordre.
\end{exercise}

\begin{exercise}[$\star\star$]\label{exo:dv:fourier-pricing-and-calibration-engineering:4}
Avec un bruit de cotations de 0.3 point de volatilité et vingt cotations, déterminer l'erreur d'ajustement qu'un modèle correct doit atteindre, et la comparer à celles du chapitre.
\end{exercise}

\begin{exercise}[$\star\star$]\label{exo:dv:fourier-pricing-and-calibration-engineering:5}
Expliquer pourquoi COS stagne à $3\times10^{-5}$ avec un intervalle de douze écarts-types, et pourquoi davantage de termes ne servent à rien.
\end{exercise}

\begin{exercise}[$\star\star$]\label{exo:dv:fourier-pricing-and-calibration-engineering:6}
Expliquer pourquoi $\kappa$ et $\eta$ évoluent ensemble le long de la vallée représentée sur la \cref{fig:dv:fourier-pricing-and-calibration-engineering:valley}.
\end{exercise}

\begin{exercise}[$\star\star\star$]\label{exo:dv:fourier-pricing-and-calibration-engineering:7}
\emph{Programmation.} Relancer la calibration sur 60 jours avec une pondération égale sur les prix au lieu de la \omterm{def:dv:fourier-pricing-and-calibration-engineering:vega}{pondération par le vega}, sans pénalité. Dire si la \omterm{def:dv:stochastic-volatility:sv}{volatilité de la
volatilité} dérive davantage ou moins, et pourquoi.
\end{exercise}

\begin{exercise}[$\star\star\star$]\label{exo:dv:fourier-pricing-and-calibration-engineering:8}
\emph{Repérer le défaut.} «~Notre calibration est excellente~: l'erreur d'ajustement est restée sous 0.45 point de volatilité chaque jour depuis deux mois.~»
\end{exercise}

\section{Problème~: le paramètre qui dérive}

\begin{problem}[{Problème du week-end --- combien pénaliser}]\label{pb:dv:fourier-pricing-and-calibration-engineering:1}
Un desk calibre chaque matin un \omterm{def:dv:stochastic-volatility:heston}{modèle de Heston} sur vingt cotations (un, trois, six et douze mois~; cinq strikes chacune) et valorise avec lui des \omterm{def:dv:variance-swaps-and-volatility-derivatives:volswap}{swaps de volatilité}
et des \omterm{def:dv:asians-lookbacks-cliquets-and-forward-starts:cliquet}{cliquets} avec lui. Sur soixante jours, la \omterm{def:dv:stochastic-volatility:sv}{volatilité de la volatilité} non pénalisée dérive.

\medskip\noindent\textbf{Partie I --- Le pricer.}
\begin{enumerate}
\item Dire lequel de COS et de Carr--Madan retenir pour cette calibration, et pourquoi.
\item Déterminer combien de termes COS donnent une erreur de prix inférieure à $10^{-4}$ sur le smile à un an.
\item Expliquer ce qui fixe l'intervalle de troncature, et ce qui se passe lorsqu'il est trop étroit.
\item Estimer grossièrement la durée d'une calibration, et ce qui la domine.
\item Indiquer comment vérifier le pricer indépendamment de la calibration.
\end{enumerate}

\medskip\noindent\textbf{Partie II --- L'ajustement.}
\begin{enumerate}[resume]
\item Indiquer ce que le filtre écarte, un jour type, sur le marché synthétique.
\item Donner les paramètres du premier jour et l'erreur d'ajustement.
\item Dire ce qu'il advient de l'ajustement à un mois avec des poids de prix égaux.
\item Évaluer la platitude de l'erreur d'ajustement le long de $\eta$, par rapport au bruit.
\item Expliquer ce que le test aller-retour avec bruit dit des paramètres que les cotations identifient.
\end{enumerate}

\medskip\noindent\textbf{Partie III --- Au jour le jour.}
\begin{enumerate}[resume]
\item Donner le plus grand mouvement quotidien de $\eta$ sans pénalité, et le jour où il se produit.
\item Dire ce qu'il fait à l'ajustement de convexité du \omterm{def:dv:variance-swaps-and-volatility-derivatives:volswap}{swap de volatilité}, et quelle est cette valeur en argent pour un \omterm{def:dv:variance-swaps-and-volatility-derivatives:veganot}{notionnel vega} de 100\,000 par point de volatilité.
\item Expliquer pourquoi ce mouvement est un \omterm{def:dv:fourier-pricing-and-calibration-engineering:recalib}{P\&L de recalibration} et non un mouvement de marché.
\item Décrire comment la pénalité modifie les trajectoires du paramètre et de l'ajustement.
\item Dire ce que coûte la pénalité en erreur d'ajustement.
\end{enumerate}

\medskip\noindent\textbf{Partie IV --- Choisir le poids.}
\begin{enumerate}[resume]
\item Décrire le rejeu qui permet de choisir $\lambda$.
\item Indiquer quel est le risque d'une pénalité trop forte.
\item Dire quels paramètres on pénaliserait davantage, et lesquels moins.
\item Énoncer le \emph{résultat nommé}~: le poids de régularisation pour lequel les variations quotidiennes de la \omterm{def:dv:stochastic-volatility:sv}{volatilité de la volatilité} passent sous 0.1 tandis que
l'erreur d'ajustement augmente de moins de 0.1 point de volatilité.
\item Résumer en une phrase ce qu'apporte une pénalité sur les variations des paramètres.
\end{enumerate}
\end{problem}

\section{Questions d'entretien}

\begin{interviewq}[$\star$ \iqroles{développeur, chercheur}]\label{iq:dv:fourier-pricing-and-calibration-engineering:1}
Comment valoriseriez-vous une option européenne lorsque vous ne disposez que de la fonction caractéristique~?
\end{interviewq}

\begin{interviewq}[$\star\star$ \iqroles{chercheur}]\label{iq:dv:fourier-pricing-and-calibration-engineering:2}
Comparez la transformée de Fourier rapide de Carr--Madan et la méthode COS.
\end{interviewq}

\begin{interviewq}[$\star\star$ \iqroles{chercheur, banque}]\label{iq:dv:fourier-pricing-and-calibration-engineering:3}
Les paramètres de votre calibration du \omterm{def:dv:stochastic-volatility:heston}{modèle de Heston} sautent d'un jour à l'autre alors que l'ajustement reste bon~: que se passe-t-il, et que faites-vous~?
\end{interviewq}

\begin{interviewq}[$\star\star$ \iqroles{développeur}]\label{iq:dv:fourier-pricing-and-calibration-engineering:4}
Concevez un traitement de calibration quotidienne. Que stocke-t-il, que vérifie-t-il et sur quoi déclenche-t-il une alerte~?
\end{interviewq}

\begin{interviewq}[$\star\star$ \iqroles{trader, risque}]\label{iq:dv:fourier-pricing-and-calibration-engineering:5}
Votre portefeuille affiche un P\&L inexpliqué important un jour calme. Comment la calibration peut-elle en être la cause~?
\end{interviewq}

\begin{interviewq}[$\star\star\star$ \iqroles{chercheur}]\label{iq:dv:fourier-pricing-and-calibration-engineering:6}
Pourquoi pondéreriez-vous une calibration, et comment choisiriez-vous les poids~?
\end{interviewq}
